\documentclass[11pt,a4paper]{article}
\usepackage[margin=1in]{geometry}
\usepackage{mathptmx}
\usepackage{booktabs}
\usepackage{longtable}
\usepackage{array}
\usepackage{xcolor}
\usepackage{enumitem}
\usepackage{titlesec}
\usepackage[colorlinks=true,linkcolor=blue!50!black,urlcolor=blue!50!black]{hyperref}
\usepackage{fancyhdr}
\usepackage{tcolorbox}

\definecolor{ink}{HTML}{1A1A1A}
\definecolor{good}{HTML}{059669}
\definecolor{warn}{HTML}{B45309}
\definecolor{bad}{HTML}{B91C1C}
\definecolor{soft}{HTML}{F4F3EE}

\newcommand{\done}{\textcolor{good}{\textbf{Done}}}
\newcommand{\partly}{\textcolor{warn}{\textbf{Partly}}}
\newcommand{\todo}{\textcolor{bad}{\textbf{Not started}}}

\titleformat{\section}{\Large\bfseries}{\thesection.}{0.5em}{}
\titleformat{\subsection}{\large\bfseries}{\thesubsection}{0.5em}{}
\titlespacing*{\section}{0pt}{18pt}{8pt}

\pagestyle{fancy}\fancyhf{}
\rhead{\small MeshSOS --- Project Report}
\lfoot{\small Generated 27 October 2026}
\rfoot{\small\thepage}
\renewcommand{\headrulewidth}{0.4pt}

\setlist[itemize]{itemsep=2pt,topsep=4pt}
\setlist[enumerate]{itemsep=2pt,topsep=4pt}

\title{\textbf{MeshSOS}\\[4pt]
\large Semantic Compression for Disaster Mesh Networks\\[10pt]
\normalsize Complete Project Report --- what is built, what is missing,\\
and what to do next}
\author{}
\date{27 October 2026}

\begin{document}
\maketitle
\thispagestyle{fancy}

\begin{tcolorbox}[colback=soft,colframe=ink!30,boxrule=0.5pt]
\textbf{How to read this.} Section 1 explains the project in plain words.
Sections 2--4 describe what exists today. Sections 5--8 cover the models,
data and results. Section 9 is the honest list of what is broken or missing
and how to fix each item. Section 10 is the schedule. If you only read one
section, read Section 9.
\end{tcolorbox}

\tableofcontents
\newpage

%====================================================================
\section{What this project is}

\subsection{The problem}

When an earthquake or flood hits, mobile towers stop working. People who need
help cannot call anyone. The phones still work --- they just have no network
to reach.

Smartphone mesh networking solves this by letting phones talk to each other
directly over Bluetooth. A message hops from phone to phone until it reaches
someone who still has internet, who forwards it to a rescue centre.

\subsection{The specific problem we work on}

Bluetooth has a hard limit. In the mode that works on \emph{every} phone ---
including cheap and old ones, which is what people actually carry in
disaster-prone areas --- one transmission carries only \textbf{31 bytes}.
After the mandatory Bluetooth headers, about \textbf{27 bytes} are usable.

Our emergency message does not fit. Measured on our own code:

\begin{center}
\begin{tabular}{lrr}
\toprule
\textbf{Message} & \textbf{Size} & \textbf{Transmissions needed} \\
\midrule
Heartbeat (``I am alive'') & 25 B & 1 \\
SOS with no text at all & 124 B & 5 \\
SOS with a short message & 161 B & 6 \\
\bottomrule
\end{tabular}
\end{center}

\textbf{The surprise:} the text is not the problem. An SOS with \emph{zero}
characters of text is still 124 bytes, because it carries a digital signature
(64 B) and a public key (32 B). Cryptography is \textbf{96 of those 124 bytes
--- 77\% of the message}.

This means compressing the text alone takes a message from 6 transmissions to
5. It can never reach 1.

\subsection{Why splitting a message is bad}

If a message needs 6 transmissions, all 6 must arrive or the message is lost.
With 3\% loss per transmission:

\begin{itemize}
\item 1 transmission: fails 3\% of the time
\item 6 transmissions: fails 16.7\% of the time
\end{itemize}

It also keeps the radio busy about six times longer, and in a crowded mesh
where every phone is relaying, that is what limits how many messages get
through.

\subsection{What we built}

A path that fits a complete, authenticated emergency report into
\textbf{one} 25-byte transmission:

\begin{enumerate}
\item The person types what happened, in English or Hindi or both.
\item The phone reads it and compresses the meaning into \textbf{3 bytes}.
\item Those 3 bytes go into a 25-byte message with location and a security tag.
\item If the phone did \emph{not} understand the message, it sends the original
      text instead --- slower, but nothing is guessed.
\end{enumerate}

%====================================================================
\section{Where the project stands today}

\begin{center}
\begin{longtable}{p{6.3cm} l p{5.5cm}}
\toprule
\textbf{Part} & \textbf{Status} & \textbf{Note} \\
\midrule
\endhead
3-byte semantic envelope & \done & Packs and unpacks, tested \\
Rule-based encoder (English + Hindi) & \done & 0 KB, 0.1 ms per message \\
Decoder and triage scoring & \done & Replaces the old hardcoded weights \\
25-byte authenticated message & \done & Fits one frame, tag survives hops \\
Size-aware mesh simulator & \done & Airtime, fragmentation, collisions \\
Network evaluation (A/B/C) & \done & 60 trials per cell \\
Stratified evaluation, 10{,}000 msgs & \done & Five strata \\
Learned model + cascade & \done & Trained, compared, reported \\
Compression baselines (gzip etc.) & \done & Answers the obvious objection \\
Android app with text input & \done & Builds; APK from GitHub Actions \\
Research paper (IEEE format) & \done & 7 pages, 4 figures \\
\midrule
Real human-written test data & \todo & Blocking a good venue \\
Measurement on a real phone & \todo & One day of work \\
Confidence intervals on results & \todo & A few hours \\
Hindi data collection & \todo & No corpus has it \\
Threat model for the security tag & \todo & One section \\
Cascade deployed in the app & \todo & Measured, not shipped \\
\bottomrule
\end{longtable}
\end{center}

%====================================================================
\section{The application, explained}

\subsection{Who the actors are}

\begin{itemize}
\item \textbf{Victim phone} --- someone in trouble sends an SOS.
\item \textbf{Relay phones} --- ordinary phones nearby that pass the message on.
      They do this in the background, without anyone touching them.
\item \textbf{Gateway phone} --- a phone that still has internet. It uploads
      everything it has heard.
\item \textbf{Rescuer dashboard} --- a web map that shows where people need help.
\end{itemize}

\subsection{What the user actually does}

\begin{enumerate}
\item Opens the app. It is already running in the background, sending a small
      ``I am alive'' heartbeat every few seconds.
\item Types what happened --- for example \emph{``3 log malbe me fanse hain,
      ek behosh''}.
\item While typing, a line below the box shows what the mesh will actually
      carry. Green means understood and compressed; orange means not
      understood, full text will be sent.
\item Presses SOS. The message goes out over Bluetooth.
\end{enumerate}

\subsection{Technologies used}

\begin{center}
\begin{tabular}{p{3.6cm} p{10.5cm}}
\toprule
\textbf{Part} & \textbf{Technology} \\
\midrule
Phone app & Kotlin, Android SDK 34, minimum Android 8 \\
Radio & Bluetooth Low Energy, legacy advertising (works on all phones) \\
Background running & Android Foreground Service with \texttt{connectedDevice} type \\
On-device NLP & Pure Kotlin rules --- \textbf{no model file, no TensorFlow, no internet} \\
Security & HMAC-SHA256 truncated to 4 bytes; Ed25519 for the full message \\
Server & Python, FastAPI, SQLite \\
Dashboard & HTML + JavaScript + Leaflet maps \\
Simulator & Python, discrete-event \\
Research models & scikit-learn (TF-IDF, logistic regression, MLP) \\
Build & GitHub Actions produces the APK automatically \\
\bottomrule
\end{tabular}
\end{center}

\subsection{How data flows, step by step}

\begin{enumerate}
\item \textbf{Text in.} The user types into the app.
\item \textbf{Encode.} \texttt{SemanticEncoder.kt} reads the text, matches
      words against its lists, and produces six values plus a confidence score.
\item \textbf{Decide.} If confidence $\geq 0.45$, compress. Otherwise send the
      original text.
\item \textbf{Pack.} \texttt{CompactSos.kt} builds a 25-byte message: header
      (13) + grid location (4) + battery (1) + envelope (3) + security tag (4).
\item \textbf{Sign.} An HMAC tag is computed over everything except the hop
      counter, so the tag still works after the message has travelled.
\item \textbf{Broadcast.} \texttt{BleTransport.kt} sends it as one Bluetooth
      advertisement.
\item \textbf{Relay.} Every phone that hears it checks whether it has seen it
      before, reduces the hop counter, waits a random moment so phones do not
      all shout at once, and re-sends.
\item \textbf{Store and forward.} If nobody is nearby, the phone keeps the
      message for up to 6 hours and re-sends every 10 seconds.
\item \textbf{Gateway.} A phone with internet unpacks the 3 bytes back into
      readable fields and uploads them.
\item \textbf{Dashboard.} The rescuer sees the location, what happened, how
      many people, and a triage score --- plus a warning if the encoder was
      unsure.
\end{enumerate}

\subsection{What the 3 bytes contain}

\begin{center}
\begin{tabular}{l r p{7.6cm}}
\toprule
\textbf{Field} & \textbf{Bits} & \textbf{Meaning} \\
\midrule
Category & 4 & trapped, collapse, medical, fire, flood, evacuation, supplies, rescue, other \\
Casualty count & 3 & 0--6 people, 7 means ``seven or more'' \\
Mobility & 2 & unknown / can walk / cannot move / unresponsive \\
Hazard flags & 6 & fire, water, gas, structural, electrical, chemical \\
Severity & 2 & unknown / none / minor / serious \\
Confidence & 3 & how sure the encoder was, 0--7 \\
Reserved & 4 & for future use \\
\bottomrule
\end{tabular}
\end{center}

Confidence is sent \emph{with the data} on purpose, so the rescuer can see
which parts of a report were guessed.

%====================================================================
\section{Existing research and where the gaps are}

\subsection{What already exists}

\begin{center}
\begin{tabular}{p{3.4cm} p{10.6cm}}
\toprule
\textbf{Work} & \textbf{What it does, and what it leaves out} \\
\midrule
HumAID, CrisisBench & Large labelled sets of disaster tweets. Solved problem.
They label what a tweet is \emph{about}, not what the emergency is. \\
\addlinespace
ResQMesh (SSRN, 2026) & Closest work. On-device classifier in a BLE mesh app.
It uses the category to \emph{order} messages. Message size never changes. \\
\addlinespace
OEPB (IETF draft, 2026) & Specifies authenticated emergency broadcast. Uses a
256-byte message and explicitly \emph{accepts} splitting on Bluetooth 4.x. \\
\addlinespace
RelayZero (blog, 2026) & Proposes the same idea as us, but with no evaluation,
no radio (a TCP socket), and numbers that contradict each other. \\
\addlinespace
Semantic communication & Active field, but assumes a channel with real
capacity. Nobody targets 27 bytes. \\
\addlinespace
CuMAC & Short security tags for small devices, built up over many messages.
Prior art for our 4-byte tag. \\
\bottomrule
\end{tabular}
\end{center}

\subsection{The gaps we found}

\begin{enumerate}
\item \textbf{Nobody targets the frame boundary.} Everyone either splits the
      message or assumes a bigger radio.
\item \textbf{Nobody measures the network effect.} Prior work reports
      classifier accuracy or queue latency, not delivery.
\item \textbf{Nobody confronts the crypto budget.} Because they do not target
      one frame, they never discover that 77\% of the message is signature.
\item \textbf{The corpora collect the wrong people.} We sampled 520 tweets
      from HumAID: 51\% commentary, 18\% news, 17\% retweets. Crisis datasets
      are written by \emph{observers}; our system carries messages from
      \emph{participants}.
\end{enumerate}

\subsection{What is new in our work}

\begin{enumerate}
\item A complete \textbf{authenticated} report in one 27-byte frame, with every
      trade stated.
\item A mesh simulator that actually models message size, and delivery numbers
      from it.
\item The measured finding that cryptography, not text, is what breaks the
      budget.
\item Two honest negative results: a synthetic-trained model that looks
      excellent and is not, and a standard corpus that does not match the
      deployment.
\end{enumerate}

%====================================================================
\section{The models}

\subsection{Model 1: the rule encoder (this is what ships)}

No training, no model file, no internet. It works by matching words against
lists.

\begin{itemize}
\item \textbf{Lexicons}: lists of words per field, in English and romanised
      Hindi together --- \texttt{trapped}, \texttt{fansa}, \texttt{malbe ke
      neeche} all live in the same list.
\item \textbf{Number extraction}: a digit counts only if a person-word appears
      within 28 characters after it, so ``4 people trapped'' is read and
      ``4th floor'' is not.
\item \textbf{Negation}: English negates \emph{before} (``no injuries''),
      Hindi negates \emph{after} (``ghayal nahi''). Both sides are checked ---
      18 characters before, 14 after.
\item \textbf{Confidence}: starts at (fields recovered)/5, then is reduced if
      most of the message was not recognised ($\times 0.6$), if the writer
      sounded unsure ($\times 0.8$), if no category matched ($\times 0.5$), or
      if the message is very short ($\times 0.7$).
\end{itemize}

\textbf{Why this matters:} if the encoder cannot read most of a message, its
confidence goes \emph{down}, not up. It treats its own confusion as a reason
not to compress.

\subsection{Model 2: the learned encoder (research only, not shipped)}

\begin{center}
\begin{tabular}{l p{9.5cm}}
\toprule
\textbf{Setting} & \textbf{Value and why} \\
\midrule
Features & TF-IDF over character n-grams, 2 to 4 characters \\
Vocabulary size & 4{,}000 (script) / 3{,}000 (notebook) \\
Weighting & Sublinear term frequency \\
Classifier (script) & Logistic regression, \texttt{max\_iter}=1000, \texttt{C}=4.0 \\
Classifier (notebook) & MLP, one hidden layer of 128 units, 30 epochs \\
Heads & One per field; one binary head per hazard bit \\
Training size & 6{,}000 generated messages \\
Training time & 9.2 s (logistic) / 5.8 s (MLP), CPU only \\
Model size & 1{,}435 KB float32, 808 KB int8 \\
Speed & 0.83 ms per message \\
\bottomrule
\end{tabular}
\end{center}

\textbf{Why character n-grams, not words:} people spell romanised Hindi
differently every time --- \emph{paani}, \emph{pani}, \emph{panee}. Character
pieces survive that; whole words do not.

\subsection{Model 3: the cascade (best result, not yet shipped)}

Run the rules first. If they are confident, use their answer. If not, pass the
message to the learned model. Nothing extra is trained --- the routing signal
is the confidence score that already existed.

%====================================================================
\section{Datasets}

\begin{center}
\begin{tabular}{p{3.4cm} r p{8.4cm}}
\toprule
\textbf{Set} & \textbf{Size} & \textbf{What it is and what it is for} \\
\midrule
Hand-written & 25 & Written by us. English and Hindi. A smoke test, not a
benchmark. \\
\addlinespace
Synthetic training & 6{,}000 & Built from templates. Labels are correct by
construction, not copied from the rule encoder. \\
\addlinespace
Stratified test & 10{,}000 & Five groups of 2{,}000: known words, synonyms the
lexicon lacks, misspellings, code-mixed, and nonsense. \\
\addlinespace
HumAID & 76{,}484 & Real tweets, downloaded. \textbf{Not usable as-is} ---
mostly commentary and news. English only. \\
\bottomrule
\end{tabular}
\end{center}

\textbf{The honest position:} there is no real human-written test data in this
project yet. That is the single biggest weakness.

\subsection{Preprocessing}

Deliberately minimal, because the encoder has to run on a phone in
milliseconds:

\begin{enumerate}
\item Unicode normalisation (NFKC) --- makes different ways of writing the
      same character identical.
\item Lowercase.
\item Collapse repeated whitespace.
\item \textbf{Nothing else.} No stemming, no stopword removal, no spell
      correction, no tokenisation into words. Matching is done on raw
      substrings so multi-word phrases like \emph{``malbe ke neeche''} work.
\end{enumerate}

%====================================================================
\section{Experimental results}

Everything below was measured on this code and can be reproduced.

\subsection{Message size}

\begin{center}
\begin{tabular}{lrr}
\toprule
 & \textbf{Bytes} & \textbf{Frames} \\
\midrule
Full SOS with text & 162--187 & 6--7 \\
Compact SOS & \textbf{25} & \textbf{1} \\
Reduction & \multicolumn{2}{r}{\textbf{about 86\%}} \\
\bottomrule
\end{tabular}
\end{center}

\subsection{Encoder accuracy by group (10{,}000 messages)}

\begin{center}
\begin{tabular}{lrrrr}
\toprule
\textbf{Group} & \textbf{Category} & \textbf{Count} & \textbf{All fields} & \textbf{Fell back} \\
\midrule
Known words      & 79.8\% & 100.0\% & 58.6\% & 16.2\% \\
Synonyms (unseen) & 8.3\% &  60.3\% &  0.1\% & 97.0\% \\
Misspellings     & 69.8\% & 100.0\% & 50.9\% & 25.2\% \\
Code-mixed       & 79.9\% & 100.0\% & 60.4\% & 17.4\% \\
Nonsense         & ---    & ---     & ---    & 100.0\% \\
\bottomrule
\end{tabular}
\end{center}

\textbf{Read the synonyms row.} When the encoder meets words it does not know,
it almost never gets the answer right --- but it refuses to compress 97\% of
the time instead of guessing. It fails \emph{slowly}, not \emph{wrongly}.

\subsection{Rules vs learned model vs cascade}

\begin{center}
\begin{tabular}{lrrr}
\toprule
\textbf{Group} & \textbf{Rules} & \textbf{Model} & \textbf{Cascade} \\
\midrule
Known words & 79.8\% & 62.3\% & \textbf{81.2\%} \\
Synonyms & 8.3\% & \textbf{39.1\%} & 38.8\% \\
Misspellings & 69.8\% & 61.7\% & \textbf{74.8\%} \\
Code-mixed & \textbf{79.9\%} & 50.5\% & 75.8\% \\
\midrule
\textbf{Overall} & 59.5\% & 53.4\% & \textbf{67.7\%} \\
\bottomrule
\end{tabular}
\end{center}

The model scored \textbf{99.7\%} on data from its own generator and
\textbf{28\%} on hand-written text. It learned the generator, not the task.
But it beats the rules by five times on unseen wording, so combining them
beats either.

\subsection{Network delivery}

60 trials per cell, 3\% loss per frame.

\begin{center}
\begin{tabular}{lrrrr}
\toprule
\textbf{Nodes} & \textbf{Full text} & \textbf{Always compressed} & \textbf{Compressed if sure} & \textbf{Airtime saved} \\
\midrule
50  & 80.0\% & 88.3\% & 86.7\% & 15.7\% \\
100 & 76.7\% & 85.0\% & 80.0\% & 12.0\% \\
200 & 53.3\% & 60.0\% & 60.0\% & 8.5\% \\
\bottomrule
\end{tabular}
\end{center}

\subsection{Why not ordinary compression}

\begin{center}
\begin{tabular}{lrrl}
\toprule
\textbf{Method} & \textbf{Average} & \textbf{Worst} & \textbf{Fits 27 B?} \\
\midrule
Raw text & 48.2 B & 65 & no \\
gzip & 65.8 B & 78 & no \\
LZMA & 105.3 B & 124 & no \\
zlib with shared dictionary & 11.2 B & 51 & sometimes \\
Our envelope & \textbf{3.0 B} & \textbf{3} & \textbf{always} \\
\bottomrule
\end{tabular}
\end{center}

Normal compressors make short messages \emph{bigger}, because their own
headers cost more than they save.

\subsection{Testing}

27 automated tests pass. One of them fails the build if anyone adds a field
that pushes the message past 27 bytes.

%====================================================================
\section{What is incomplete, and how to fix it}

This is the important section. Each item says what is wrong, why it matters,
and exactly how to fix it.

\subsection{No real human-written test data}

\textbf{Status:} \todo. \textbf{This blocks publication at a good venue.}

\textbf{Why it matters.} Every accuracy number in this project comes from text
that we or a generator wrote. A reviewer will say the evaluation measures our
own writing style.

\textbf{How to fix.}
\begin{enumerate}
\item Run \texttt{notebooks/MeshSOS\_build\_annotation\_set.ipynb} in Colab.
      It downloads HumAID and writes two labelling sheets.
\item \textbf{Filter first.} The notebook now counts how many messages are
      genuine first-person distress. If that number is small, use the
      \emph{Disaster Response Messages} corpus instead --- its
      \texttt{direct} genre is real SMS sent during the Haiti and Chile
      earthquakes.
\item Two people label 300--500 messages \textbf{independently}, following
      \texttt{docs/ANNOTATION\_GUIDE.md}. Not together --- agreement measured
      after discussion means nothing.
\item Run \texttt{python nlp/score\_annotations.py A.csv B.csv}. It reports
      Cohen's kappa per field and lists disagreements.
\item Settle the disagreements, build \texttt{gold.csv}, then run the same
      script with \texttt{--evaluate}.
\end{enumerate}

\textbf{Effort:} 6--8 hours per person, two people.

\subsection{Nothing has run on a real phone}

\textbf{Status:} \todo. \textbf{Second biggest weakness.}

\textbf{Why it matters.} All latency figures are desktop Python. There is no
battery measurement. For a paper about phones, that is a visible hole.

\textbf{How to fix.}
\begin{enumerate}
\item Install the APK (already built by GitHub Actions) on one cheap Android
      phone, four years old or more.
\item Add timing around the \texttt{encode()} call --- use
      \texttt{System.nanoTime()} and log it.
\item Type 50 messages, record the times, report mean and worst case.
\item Note the APK size difference with and without the encoder.
\item If time allows, use Android's Battery Historian for drain over an hour.
\end{enumerate}

\textbf{Effort:} one day. \textbf{This is the cheapest big improvement.}

\subsection{No confidence intervals on the network results}

\textbf{Status:} \todo.

\textbf{Why it matters.} The 100-node result is $+4.3\%$. With 60 trials and no
error bars, nobody can tell whether that is real or noise. A reviewer will
write exactly that.

\textbf{How to fix.} In \texttt{simulator/eval\_compression.py}, record each
trial's outcome instead of only the total, then report mean $\pm$ 95\%
confidence interval, and run a two-proportion test between A and C. Add error
bars to the delivery figure.

\textbf{Effort:} 2--3 hours.

\subsection{The Hindi claim has no Hindi data}

\textbf{Status:} \todo.

\textbf{Why it matters.} The paper claims code-mixed support and reports
60.4\% on it --- but that is on messages we generated. HumAID is English only.
No public disaster corpus has Hindi.

\textbf{How to fix.} Give 25--30 classmates a scenario --- \emph{``imagine you
are trapped after an earthquake; write the message you would send''} --- and
collect 5 each. That is 150 messages in a day, written by people who are not
the authors. Describe it in the paper as \emph{elicited}, not collected in the
wild. Honest and acceptable.

\textbf{Effort:} one day.

\subsection{The security tag is weak and has no threat model}

\textbf{Status:} \partly{} --- the tag works, the analysis does not exist.

\textbf{Why it matters.} A 4-byte tag means a 1-in-4-billion forgery chance
per attempt. The paper says so but never says who the attacker is or what they
can do.

\textbf{How to fix.}
\begin{enumerate}
\item Write a short threat-model subsection: who the attacker is, what they can
      send, what they cannot do, and what the gateway does about rate limits.
\item \textbf{Better:} implement the CuMAC idea. Split a full 16-byte tag into
      four 4-byte pieces and send a different piece in each re-broadcast. The
      store-and-forward layer already repeats every 10 seconds for 6 hours, so
      the repeats are free. The gateway acts on the first frame and reaches
      full strength over later ones.
\end{enumerate}

\textbf{Effort:} 2 hours for the section, 2 days for the implementation.

\subsection{The cascade is measured but not deployed}

\textbf{Status:} \partly{}.

\textbf{Why it matters.} The best result in the paper (67.7\%) is not what the
app does. The app runs rules only.

\textbf{How to fix.} Either ship it --- which means putting a quantised model
in the APK and adding TensorFlow Lite --- or state clearly in the paper that
it is an offline result. We currently do the second, which is honest. Shipping
it is a real engineering job because it adds a model file and a dependency.

\textbf{Effort:} 3--4 days to ship properly.

\subsection{The fallback path is wasteful}

\textbf{Status:} \todo{} --- a known easy improvement.

\textbf{Why it matters.} When the encoder is unsure (23\% of messages), it
sends raw text --- 162--187 bytes, 6--7 frames. Our own measurements show
shared-dictionary compression would bring that to about 11 bytes.

\textbf{How to fix.} Build a dictionary from common emergency phrases, ship it
with the app, and compress the fallback text with \texttt{zlib} using that
dictionary. Both ends must have the same dictionary.

\textbf{Effort:} one day.

\subsection{The radio model is not validated}

\textbf{Status:} \partly{}.

\textbf{Why it matters.} We wrote the collision and airtime model ourselves. It
has not been checked against a real measurement or a second simulator.

\textbf{How to fix.} Either compare against a known BLE simulator, or measure
real airtime on two phones with a sniffer, or state it clearly as a limitation
(currently done).

\textbf{Effort:} days to weeks. Lowest priority of the list.

\subsection{Old results in \texttt{docs/EVALUATION.md} are single-seed}

\textbf{Status:} \partly{} --- known and documented.

\textbf{Why it matters.} The original 200-node ``0 delivered'' figure came from
one unlucky random seed. The median across seeds is 6 hops, not 17.

\textbf{How to fix.} Re-run \texttt{simulator/evaluation\_matrix.py} over 25
seeds and regenerate the file with mean and standard deviation.

\textbf{Effort:} one hour.

%====================================================================
\section{Roadmap}

\subsection{Week 1 --- make the numbers defensible}

\begin{center}
\begin{tabular}{p{1.6cm} p{8.4cm} p{3.4cm}}
\toprule
\textbf{Day} & \textbf{Task} & \textbf{Output} \\
\midrule
1 & Confidence intervals and significance tests on all network results & Updated tables, error bars \\
1 & Re-run old evaluation over 25 seeds & New \texttt{EVALUATION.md} \\
2 & Install APK on one old phone; time the encoder; note APK size & Real latency numbers \\
2 & Write the threat-model subsection & New \S{} in the paper \\
3--5 & Collect 150 Hindi messages from classmates & First real data \\
\bottomrule
\end{tabular}
\end{center}

\subsection{Weeks 2--3 --- real evaluation data}

\begin{enumerate}
\item Decide the corpus: HumAID filtered, or Disaster Response Messages
      \texttt{direct} genre. Read samples from both before choosing.
\item Build the annotation sheets in Colab.
\item Two people label 300--500 messages independently.
\item Score agreement; settle disagreements; build the gold set.
\item Re-run every experiment against it. The scripts already exist --- only
      the data changes.
\end{enumerate}

\subsection{Week 4 --- improvements worth making}

\begin{enumerate}
\item Compress the fallback path with a shared dictionary.
\item Implement segmented authentication over the store-and-forward repeats.
\item Optionally ship the cascade in the app.
\end{enumerate}

\subsection{Week 5 --- paper}

\begin{enumerate}
\item Replace every synthetic number with the real one.
\item Add the hardware measurements.
\item Add error bars to all figures.
\item Re-read the Related Work; check whether the FMEC 2026 ResQMesh paper has
      appeared.
\item Have someone outside the team read it for clarity.
\end{enumerate}

\subsection{Honest assessment of where it can go}

\begin{center}
\begin{tabular}{p{6.4cm} p{3.0cm} p{4.4cm}}
\toprule
\textbf{Venue} & \textbf{Today} & \textbf{After the roadmap} \\
\midrule
College / university symposium & Strong & Strong \\
National conference & Good & Strong \\
IEEE / ACM workshop & Possible & Good \\
Mid-tier conference & Unlikely & Possible \\
Top venue (SenSys, MobiCom) & No & Still unlikely \\
\bottomrule
\end{tabular}
\end{center}

The two items that move this the most are \textbf{real data} and \textbf{one
phone measurement}. Everything else is polish.

%====================================================================
\section{How to run everything}

\begin{verbatim}
pip install -r requirements.txt

python demo.py                                 # end-to-end walkthrough
python -m pytest tests/ -q                     # 27 tests
python simulator/test_phase0_size_aware.py     # size-awareness check
python nlp/eval_encoders.py 2000               # 10,000-message evaluation
python nlp/train_model.py                      # the learned model
python simulator/eval_compression.py 60        # network results
python simulator/make_figure.py                # delivery chart
python paper/make_figures.py                   # all paper figures
python simulator/ttl_vs_size_experiment.py 25  # TTL vs size
\end{verbatim}

Backend and dashboard: \texttt{./start.sh} or \texttt{docker compose up}.

The Android APK is built automatically by GitHub Actions on every push;
download it from the workflow run's artifacts.

\vspace{1em}
\begin{tcolorbox}[colback=soft,colframe=ink!30,boxrule=0.5pt]
\textbf{One last thing.} Every number in this report came from running the
code in this repository. If anyone asks where a figure came from, the command
that produced it is listed above. That is worth more than any single result.
\end{tcolorbox}

\end{document}
